\section{模型的评价、改进与推广}

\subsection{模型优点}

本文的首要优点是将四个问题组织为可验收的分层证据链。问题一只在附件 A 内建立质量与配比响应，问题二只在附件 B 内估计绝对 Loss 标度律，跨体系信息通过 $Q_{\mathrm{ref}}$、$\boldsymbol p_{\mathrm{ref}}$ 和 $R(\boldsymbol p)$ 等显式接口传递，避免了将实验对象不同的 Loss 直接拼接。每一次向后传递都附带验收状态和适用边界，使“已识别的中心模型”与“仅用于敏感性的情景扩展”保持区分。

其次，各子模型均保留了对应的结构约束。质量评价对指标进行语义同向化并保留块间冲突；配比响应在单纯形上使用 Scheff\'e 混料结构；标度律对系数施加正值和单调性约束；资源配置同时考虑预算、质量上下界和观测支持域。问题三还利用目标对 $D$ 的单调性将三维优化解析降为二维，并以完整三维优化、加密网格和多起点轨迹交叉复核，降低了局部最优和单一初值造成的风险。

再次，模型不仅给出点估计，还通过分块交叉验证、外部对照、bootstrap、参数扰动和支持域标记报告不确定性。对于质量成本、质量尺度和长上下文等无法由当前附件唯一识别的要素，本文使用多情景比较而非将某一假设包装为确定事实，从而使结论的证据强度与表述强度相匹配。

\subsection{模型局限}

问题一中 $11$ 个推断域缺少与七个质量锚点等强的直接对齐证据，因而其领域质量回退到公共参考水平；这保证了稳健性，却牺牲了细粒度区分能力。配比响应面在 $1\mathrm{M}$、$60\mathrm{M}$ 和 $1\mathrm{B}$ 实测尺度上具有较好排序能力，但近优配比明显受支持半径、份额上界和局部解簇影响，不能直接解释为大模型训练中的唯一最优数据配方。

问题二的主要局限是跨体系可识别性。质量尺度只能以 $Q_{\mathrm{ref}}$ 为锚点作情景映射，配比响应 $R(\boldsymbol p)$ 则缺少映射到 B 体系绝对 Loss 的强度校准。因此，问题三必须固定 $\boldsymbol p_{\mathrm{ref}}$，并不能完整回答“配比与规模联合自由优化”的强形式问题。同时，高预算配置超出 B1/B7 的共同观测区间，三类质量成本又是题设情景而非实测工程成本，故问题三的高预算解应被视为条件外推。

问题四的确认性能力前沿样本只覆盖 $7$ 个月和 $13$ 个模型家族，非规模残余中还混合了数据工程、架构、优化和后训练等未充分观测因素。因此，当前贡献分解属于观察性分解，$12$ 个月与 $24$ 个月的预测属于带宽区间的情景外推，不应被解释为确定性或严格因果结论。

\subsection{改进方向}

最直接的改进是补充跨体系共同实验。可在同一模型族、同一验证集和统一训练流程下，对 $(N,D,Q_A,\boldsymbol p)$ 实施分层因子实验，为配比迁移强度 $\lambda_0$ 和质量尺度斜率 $s_Q$ 提供可识别数据。当共同实验覆盖充足时，可将固定配比优化扩展为配比、质量和规模的联合优化，并直接检验二阶交互是否随模型规模保持。

对资源配置模型，可进一步收集语料过滤、去重、合成、审核和存储的实际成本，将情景函数 $g(Q_A)$ 替换为经验成本曲线；同时把通信、内存、混合精度和硬件利用率纳入计算成本，将 FLOPs 约束扩展为多资源约束。对超出当前支持域的配置，应增加实际训练锚点或使用保守鲁棒优化，而不仅依赖幂律外推。

对能力前沿预测，应延长同口径时间面板，增加架构、数据处理、优化器、后训练算力和许可证等可比特征，并在更多历史起点上进行滚动回测。这将有助于把当前的“非规模残余”进一步拆分为更可解释的技术通道，并降低长期外推区间对短面板状态斜率的依赖。

\subsection{模型推广}

本文的分层框架不依赖某一特定大语言模型。对视觉基础模型、多模态模型或科学计算模型，只需把 Token 数替换为对应的样本、图像块或模态单元，重新定义质量信号和长序列成本，即可沿用“数据质量——性能标度——资源优化”的建模路径。在实际训练管理中，该框架也可作为情景规划工具：用可观测区间内的中心解给出基准配置，用质量成本、上下文和外推边界情景刻画决策风险，而不将单一点估计作为无条件的工程承诺。

\bibliographystyle{gmcm}
\bibliography{reference}
